% !TEX root = MinImageScorer_GCCE2026_v3.tex
\documentclass[conference]{IEEEtran}
\IEEEoverridecommandlockouts

\usepackage{cite}
\usepackage{amsmath,amssymb,amsfonts}
\usepackage{graphicx}
\usepackage{textcomp}
\usepackage{xcolor}
\usepackage{booktabs}
\usepackage{multirow}
\usepackage{microtype}

\def\BibTeX{{\rm B\kern-.05em{\sc i\kern-.025em b}\kern-.08em
    T\kern-.1667em\lower.7ex\hbox{E}\kern-.125emX}}

\begin{document}

\title{MinImageScorer: Error-Driven Difficulty Scoring\\
for Guided Copy-Paste Augmentation in\\
Agricultural Object Detection}

% TODO: Fill in name, supervisor, affiliation
\author{\IEEEauthorblockN{1\textsuperscript{st} Given Name Surname}
\IEEEauthorblockA{\textit{dept. name of organization (of Aff.)} \\
\textit{name of organization (of Aff.)}\\
City, Country \\
email address or ORCID}
\and
\IEEEauthorblockN{2\textsuperscript{nd} Given Name Surname}
\IEEEauthorblockA{\textit{dept. name of organization (of Aff.)} \\
\textit{name of organization (of Aff.)}\\
City, Country \\
email address or ORCID}
\and
\IEEEauthorblockN{3\textsuperscript{rd} Given Name Surname}
\IEEEauthorblockA{\textit{dept. name of organization (of Aff.)} \\
\textit{name of organization (of Aff.)}\\
City, Country \\
email address or ORCID}
}

\maketitle

% ----------------------------------------------------------
\begin{abstract}
We propose \textbf{MinImageScorer}, an error-driven framework that
scores per-object detection difficulty from baseline model failures,
combining confidence and localization penalties calibrated via
Pearson correlation optimization. Validated on MinneApple and Global
Wheat Head (GWHD), the score reliably ranks objects and images by
difficulty (Pearson $r$ up to $0.82$ with miss rate) and recovers
domain- and size-level failure patterns without supervision.
Score-guided copy-paste outperforms random selection on MinneApple
($\Delta\mathrm{AP}{=}{+}0.013$) but both strategies fall below
baseline due to distributional shift. On GWHD, domain shift renders
copy-paste ineffective under either strategy. These results show
that MinImageScorer correctly identifies hard objects, and that
copy-paste utility is determined by the dataset's failure mode, not
the selection strategy.
\end{abstract}

\begin{IEEEkeywords}
object detection, copy-paste augmentation, difficulty scoring,
agricultural detection, small objects, domain shift
\end{IEEEkeywords}

% ----------------------------------------------------------
\section{Introduction}

Agricultural object detection benchmarks such as
MinneApple~\cite{hani2020minneapple} and Global Wheat Head
(GWHD)~\cite{david2021global} exhibit distinct structural failure
modes: MinneApple is dominated by sub-resolution small objects, while
GWHD has severe cross-domain visual variation. Copy-paste augmentation
has only demonstrated gains on MS COCO~\cite{kisantal2019augmentation,
ghiasi2021simple}, but do not showcase it on differnt closer domain dataset. 
Applying copy and paste augmentation blindly to structurally different datasets 
risks corrupting the training distribution rather than improving it, yet no prior work
provides an empirical framework to diagnose this mismatch.

We address this gap with three contributions: (1)~\textbf{MIS}
(\textbf{M}in\textbf{I}mage\textbf{S}corer), a per-object difficulty
score derived from model failure signals with data-driven weight
calibration; (2)~empirical validation showing the score correctly
recover *(this part is a little bit off) image-, domain-, and object-size-level difficulty without
supervision; and (3)~a controlled study showing score-guided
copy-paste outperforms random selection on a resolution-limited dataset,
while both strategies fail on a domain-shifted dataset.

% ----------------------------------------------------------
\section{Methodology}

\subsection{Scoring Formulation}

Let $G(I)$ be ground-truth objects in image $I$ and
$\mathcal{P}(g) = \{p : \mathrm{IoU}(p,g) \ge \tau\}$
the matched predictions. Denoting $c_p = \mathrm{conf}(p)$
and $u_p = \mathrm{IoU}(p,g)$, the \emph{object difficulty score} is:

\begin{equation}
S_{\mathrm{obj}}(g) =
\begin{cases}
  \min_{p \in \mathcal{P}(g)}
    \bigl[\alpha(1{-}c_p) + \beta(1{-}u_p)\bigr],
    & \mathcal{P}(g) \neq \varnothing, \\[3pt]
  \alpha + \beta, & \text{otherwise,}
\end{cases}
\label{eq:obj}
\end{equation}

where $\alpha$ and $\beta$ weight confidence and localization penalties.
Missed objects receive the maximum score $\alpha{+}\beta$. The
\emph{image-level score} combines mean object difficulty with
false-positive rate:

\begin{equation}
S_{\mathrm{img}}(I)
  = w_1 \cdot \frac{1}{|G(I)|}\sum_{g\in G(I)} S_{\mathrm{obj}}(g)
  \;+\; w_2 \cdot \mathrm{FPrate}(I).
\label{eq:img}
\end{equation}

Object scoring uses low-confidence predictions ($\tau{=}0.01$) to
capture near-misses; FP rate uses optimal-threshold predictions
to reflect deployment errors. Weights $\alpha, \beta, w_1, w_2$
are calibrated by maximizing Pearson correlation with per-image
miss rate and FP rate on the training set, requiring no
domain-specific assumptions.

\subsection{Score-Guided Copy-Paste Pipeline}

A baseline detector *(in this i only train with yolov8) 
is trained . Inference on the training set yields confidence and IoU
signals used to compute $S_{\mathrm{obj}}$ and $S_{\mathrm{img}}$.
For score-guided augmentation, objects are sampled proportionally
to $S_{\mathrm{obj}}$ via exponential weighting; for random, sampling
is uniform. Both conditions paste 2--8 objects per image and retrain
on the augmented set under identical hyperparameters.
*(This part i think we should disscuss more about the copy-paste method, like how to paste, 
how to avoid occlusion, etc. I can add more details if we decide to keep it.)

% ----------------------------------------------------------
\section{Score Validation}
\label{sec:validation}

\subsection{Correlation with Detection Failures}

Table~\ref{tab:corr} reports Pearson correlation between
$S_{\mathrm{img}}$ and per-image miss rate and FP rate on training
images. All cases show strong positive correlation ($r \ge 0.69$).
GWHD with traditional augmentation reaches $r{=}0.82$ on both
metrics, where the baseline model is stable and the score signal is
clean. On MinneApple without traditional augmentation, inter-seed
stability drops to $r{\approx}0.67$: without color jitter and
flipping, unstabilized models produce fluctuating confidence on
densely clustered small objects, making the score training-run
dependent. *(This data was not in the table which might make it hard for user to udnerstand) 
*(Also say what has been prove from this table. like what is this experiemt imply 
what is the score actually capture, etc.)

\begin{table}[t]
\centering
\caption{Pearson correlation of $S_{\mathrm{img}}$ with per-image
miss rate and FP rate (training set, w/ traditional augmentation).}
\label{tab:corr}
\resizebox{\columnwidth}{!}{%
\begin{tabular}{@{}lcccc@{}}
\toprule
\textbf{Dataset}
  & $r$(\,score,\,miss rate\,)
  & $r$(\,score,\,FP rate\,)
  & Inter-seed $r$ \\
\midrule
MinneApple  & 0.698 & 0.688 & 0.806 \\
GWHD        & 0.816 & 0.818 & 0.900 \\
\bottomrule
\end{tabular}%
}
\end{table}

\subsection{Domain- and Size-Level Validity}

Beyond per-image correlation, the score recovers structural
difficulty without access to test labels. On GWHD, mean
$S_{\mathrm{img}}$ inversely tracks AP across all seven source
domains (Spearman $\rho = -0.93$): the hardest domain
(\textit{arvalis~2}, AP\,=\,0.365) scores 0.431, while the
easiest (\textit{rres~1}, AP\,=\,0.572) scores only 0.260.
On MinneApple, small objects (${\le}32$\,px) average a score
of 0.378 vs.\ 0.222 for medium objects, exactly matching the AP
gap (0.132 vs.\ 0.522). (*I think this part is make the reader hard to imageine
we should add a table or figure to show the domain-level and size-level score distribution, which can directly show the relationship between the score and the difficulty at different level. I can add it if we decide to keep this part.)

These results confirm that MIS captures
image-, domain-, and object-size-level difficulty from training-set
model failures alone.

% ----------------------------------------------------------
\section{Augmentation Experiments}

\subsection{Setup}

All models use YOLOv8s with COCO pretrained weights.
%
% [FILL] Insert exact dataset split counts. Example:
%   MinneApple: 536 train / 67 val / 67 test images.
%   GWHD: 2,700 train / 735 val / 748 test images.
%
Training: MinneApple 200 epochs / patience 10; GWHD 100 epochs /
patience 8; batch 16; input $1024{\times}1024$; lr $10^{-2}$.
All built-in augmentations are disabled to isolate copy-paste.
Augmentation ratio 0.3; scale jitter $[0.9,1.1]$ (MA) and
$[0.8,1.2]$ (GWHD). 
*(Do we need to write all of this information ?)

\subsection{Results}

Table~\ref{tab:ap} shows that on MinneApple, score-guided copy-paste
outperforms random across all sub-metrics ($\Delta\mathrm{AP}
{=}{+}0.013$), confirming the score directs augmentation toward
higher-signal objects. However, both conditions fall below baseline
($-0.024$ for Score CP). The AP$_{75}$ drop ($0.452{\to}0.440$)
indicates degraded localization precision from pasting hard objects
without blending or masks; the AP$_S$ drop ($0.300{\to}0.286$)
indicates score-selected objects are predominantly sub-resolution
and cannot be made more learnable by repetition.

On GWHD, score-guided and random selection are indistinguishable
($\Delta\mathrm{AP}{=}-0.002$) and both fall below baseline ($-0.023$).
Score-guided sampling concentrates objects from domain-minority
sources (\textit{arvalis}), amplifying rather than correcting the
existing domain imbalance. Backbone feature analysis (YOLOv8s P4,
UMAP) confirms that augmented images are displaced from the test
distribution relative to original training images on both datasets,
with score-guided images showing greater displacement on GWHD---
directly explaining the consistent degradation.

% ----------------------------------------------------------------
% FIGURE RECOMMENDATION: produce a scatter plot of S_img vs. per-image
% miss rate for MA (left panel) and GWHD (right panel). One point per
% training image; regression line shown. This is the direct visual
% proof of Table I and takes ~4cm of column height.
% Export as score_validation.pdf and replace the block below.
% ----------------------------------------------------------------
\begin{figure}[t]
\centering
%\includegraphics[width=\columnwidth]{score_validation.pdf}
\framebox{\parbox{0.88\columnwidth}{\centering\vspace{1.2cm}
\small REPLACE with scatter: $S_{\mathrm{img}}$ vs.\ miss rate\\
Left: MinneApple ($r{=}0.70$) \quad Right: GWHD ($r{=}0.82$)\\
Each point = one training image. Regression line shown.
\vspace{1.2cm}}}
\caption{$S_{\mathrm{img}}$ vs.\ per-image miss rate.
  Higher score predicts harder images on both datasets.
  Points generated from training-set inference only.}
\label{fig:scatter}
\end{figure}

\begin{table}[t]
\centering
\caption{Three-seed mean COCO AP. \textbf{Bold} = best
  copy-paste condition per dataset.}
\label{tab:ap}
\resizebox{\columnwidth}{!}{%
\begin{tabular}{@{}llrrrrr@{}}
\toprule
\textbf{Dataset} & \textbf{Condition}
  & \textbf{AP}
  & \textbf{AP\textsubscript{50}}
  & \textbf{AP\textsubscript{75}}
  & \textbf{AP\textsubscript{S}}
  & \textbf{AP\textsubscript{M}} \\
\midrule
\multirow{3}{*}{MinneApple}
  & Baseline          & 0.439 & 0.771 & 0.452 & 0.300 & 0.592 \\
  & Random CP         & 0.415 & 0.741 & 0.427 & 0.279 & 0.563 \\
  & \textbf{Score CP} & \textbf{0.428} & \textbf{0.759}
                      & \textbf{0.440} & \textbf{0.286} & \textbf{0.581} \\
\midrule
\multirow{3}{*}{GWHD}
  & Baseline          & 0.507 & 0.922 & 0.491 & 0.136 & 0.501 \\
  & Random CP         & 0.486 & 0.905 & 0.456 & 0.118 & 0.480 \\
  & \textbf{Score CP} & \textbf{0.484} & \textbf{0.905}
                      & \textbf{0.455} & \textbf{0.098} & \textbf{0.477} \\
\bottomrule
\end{tabular}%
}
\end{table}

% ----------------------------------------------------------
\section{Conclusion}

MIS (*Say the full name) reliably quantifies detection difficulty from model failure
signals, validated by per-image correlations ($r$ up to $0.82$),
near-perfect domain difficulty ranking (Spearman $\rho{=}-0.93$
on GWHD), and correct object-size ordering on MinneApple. 
Score-guided copy-paste outperforms random
selection on MinneApple ($+0.013$~AP) but both strategies degrade
below baseline due to distributional shift between augmented and
test images. On GWHD, domain shift makes copy-paste ineffective
regardless of selection quality. The practical implication is
direct: copy-paste is effective for object underrepresentation
(COCO-like settings) but not for resolution-limited or
domain-shifted datasets; MIS provides the diagnostic to determine
which condition holds before augmentation is applied. Future work
includes domain-constrained pasting and online augmentation to
reduce distributional rigidity. (*Say this more clearly, 
like what is the future direction, what is the potential solution, etc.)

% ----------------------------------------------------------
\begin{thebibliography}{00}

\bibitem{hani2020minneapple}
N.~Hani, P.~Roy, and B.~Bhanu,
``MinneApple: A Benchmark Dataset for Apple Detection and
Segmentation,''
\textit{IEEE Robotics and Automation Letters},
vol.~5, no.~2, pp.~852--858, 2020.

\bibitem{david2021global}
E.~David \textit{et~al.},
``Global Wheat Head Detection (GWHD) Dataset,''
\textit{Plant Phenomics}, vol.~2021, 2021.

\bibitem{kisantal2019augmentation}
M.~Kisantal \textit{et~al.},
``Augmentation for Small Object Detection,''
in \textit{Proc.\ ICACCI}, 2019.

\bibitem{ghiasi2021simple}
G.~Ghiasi \textit{et~al.},
``Simple Copy-Paste is a Strong Data Augmentation Method
for Instance Segmentation,''
in \textit{Proc.\ IEEE/CVF CVPR}, pp.~2918--2928, 2021.

\end{thebibliography}

\end{document}
